% Source: https://artofproblemsolving.com/community/c3h2271569_2001_state_sprint_problem_24
% Stable identity: aops-topic-2271569-q01
\textbf{(a)} List the partitions by their largest part:
\begin{itemize}
\item largest part $7$: $7$;
\item $6$: $6+1$;
\item $5$: $5+2$, $5+1+1$;
\item $4$: $4+3$, $4+2+1$, $4+1+1+1$;
\item $3$: $3+3+1$, $3+2+2$, $3+2+1+1$, $3+1+1+1+1$;
\item $2$: $2+2+2+1$, $2+2+1+1+1$, $2+1+1+1+1+1$;
\item $1$: $1+1+1+1+1+1+1$.
\end{itemize}
That is $1+1+2+3+4+3+1=\boxed{15}$ partitions.

\textbf{(b)} Write $7$ as seven $1$s in a row. An ordered sum is obtained by choosing which of the $6$ gaps between neighbouring $1$s get a bar (stars and bars), so there are $2^6=\boxed{64}$ ordered sums. Stars and bars distinguishes the order of the parts: the partition $4+2+1$ alone corresponds to $3!=6$ ordered sums, while $3+3+1$ corresponds to only $3$. Because different partitions correspond to different numbers of ordered sums, there is no simple division that turns $64$ into the number of partitions; partitions are counted by casework (or by recursion), not by stars and bars. (For $4$: $2^3=8$ ordered sums but $5$ partitions.)
